% PHẦN 1: Nội dung phía trên khung Quy tắc thương
\noindent
\begin{minipage}[t]{0.24\textwidth}
    \mbox{}
\end{minipage}%
\hfill
\begin{minipage}[t]{0.73\textwidth}
    \fontsize{10}{14}\selectfont
    Ví dụ 2 cho thấy rằng đôi khi việc rút gọn một tích các hàm số trước khi lấy đạo hàm sẽ dễ dàng hơn so với việc sử dụng Quy tắc nhân. Tuy nhiên, trong Ví dụ 1, Quy tắc nhân là phương pháp khả dĩ duy nhất.

    \vspace{8pt}
    \noindent\textbf{\textsf{\color{stewartcyan}VÍ DỤ 3}}\hspace{0.8em}Nếu $f(x) = \sqrt{x}\,g(x)$, trong đó $g(4) = 2$ và $g'(4) = 3$, hãy tìm $f'(4)$.

    \vspace{6pt}
    \noindent\textbf{\textsf{\color{stewartcyan}LỜI GIẢI}}\hspace{0.8em}Áp dụng Quy tắc nhân, ta được
    \begin{align*}
    f'(x) &= \frac{d}{dx}\left[\sqrt{x}\,g(x)\right] = \sqrt{x}\,\frac{d}{dx}[g(x)] + g(x)\,\frac{d}{dx}[\sqrt{x}] \\[4pt]
    &= \sqrt{x}\,g'(x) + g(x)\cdot \tfrac{1}{2}x^{-1/2} \\[4pt]
    &= \sqrt{x}\,g'(x) + \frac{g(x)}{2\sqrt{x}}
    \end{align*}
    \vspace{-4pt}
    \noindent Do đó\hspace{2em}$\displaystyle f'(4) = \sqrt{4}\,g'(4) + \frac{g(4)}{2\sqrt{4}} = 2\cdot 3 + \frac{2}{2\cdot 2} = 6.5$\hfill{\color{stewartcyan}\rule{1.15ex}{1.15ex}}

    \vspace{10pt}
    \noindent{\color{stewartcyan}\rule{1.15ex}{1.15ex}}\hspace{0.5em}\textbf{\large Quy tắc thương}

    \vspace{5pt}
    Chúng ta tìm quy tắc tính đạo hàm của thương hai hàm khả vi $u = f(x)$ và $v = g(x)$ theo cách hoàn toàn tương tự như cách chúng ta đã tìm Quy tắc nhân. Nếu $x, u,$ và $v$ thay đổi các lượng tương ứng là $\Delta x, \Delta u,$ và $\Delta v$, thì lượng thay đổi tương ứng của thương $u/v$ là
    \begin{align*}
    \Delta\left(\frac{u}{v}\right) &= \frac{u + \Delta u}{v + \Delta v} - \frac{u}{v} = \frac{(u + \Delta u)v - u(v + \Delta v)}{v(v + \Delta v)} \\[6pt]
    &= \frac{v\,\Delta u - u\,\Delta v}{v(v + \Delta v)}
    \end{align*}
    do đó
    \[
    \frac{d}{dx}\left(\frac{u}{v}\right) = \lim_{\Delta x \to 0} \frac{\Delta(u/v)}{\Delta x} = \lim_{\Delta x \to 0} \frac{v\,\dfrac{\Delta u}{\Delta x} - u\,\dfrac{\Delta v}{\Delta x}}{v(v + \Delta v)}
    \]
    Khi $\Delta x \to 0$, thì $\Delta v \to 0$ vì $v = g(x)$ là hàm khả vi và do đó liên tục. Như vậy, sử dụng các Quy tắc giới hạn, ta có
    \[
    \frac{d}{dx}\left(\frac{u}{v}\right) = \frac{v \displaystyle\lim_{\Delta x \to 0}\frac{\Delta u}{\Delta x} - u \displaystyle\lim_{\Delta x \to 0}\frac{\Delta v}{\Delta x}}{v \displaystyle\lim_{\Delta x \to 0}(v + \Delta v)} = \frac{v\,\dfrac{du}{dx} - u\,\dfrac{dv}{dx}}{v^2}
    \]
\end{minipage}

\vspace{8pt}
% PHẦN 2: Khung Quy tắc thương dóng ngang hàng 1:1 với ghi chú lề trái
\noindent
\begin{minipage}[t]{0.24\textwidth}
    \vspace{2pt}
    \fontsize{8.5}{11.5}\selectfont
    \raggedright
    Dưới dạng ký hiệu dấu phẩy, Quy tắc thương được viết là
    \vspace{1pt}
    \[
    \left(\frac{f}{g}\right)' = \frac{g f' - f g'}{g^2}
    \]
\end{minipage}%
\hfill
\begin{minipage}[t]{0.73\textwidth}
    \vspace{0pt}
    \begin{definitionbox}
    \textbf{\color{stewartred}Quy tắc thương}\hspace{1em}Nếu $f$ và $g$ là các hàm khả vi, thì
    \[
    \frac{d}{dx}\left[\frac{f(x)}{g(x)}\right] = \frac{g(x)\,\dfrac{d}{dx}[f(x)] - f(x)\,\dfrac{d}{dx}[g(x)]}{[g(x)]^2}
    \]
    \end{definitionbox}
\end{minipage}

\vspace{8pt}
% PHẦN 3: Diễn giải bằng lời dưới khung
\noindent
\begin{minipage}[t]{0.24\textwidth}
    \mbox{}
\end{minipage}%
\hfill
\begin{minipage}[t]{0.73\textwidth}
    \fontsize{10}{14}\selectfont
    Bằng lời, Quy tắc thương phát biểu rằng \textit{đạo hàm của một thương bằng mẫu số nhân với đạo hàm của tử số trừ đi tử số nhân với đạo hàm của mẫu số, tất cả chia cho bình phương của mẫu số.}
\end{minipage}
